\documentclass[11pt]{article}
\usepackage[utf8]{inputenc}
\usepackage[T1]{fontenc}
\usepackage{geometry}
\usepackage{hyperref}
\geometry{margin=1in}

\title{\textbf{JACM Submission Cover Letter}}
\author{}
\date{}

\begin{document}

\maketitle

\noindent\textbf{Date}: March 31, 2026

\noindent\textbf{To}: Professor Venkatesan Guruswami, Editor-in-Chief \\
\textbf{Journal of the ACM}

\noindent\textbf{Subject}: Submission of ``The Optimizer Quotient and the Certification Trilemma''

Dear Professor Guruswami,

Please accept our submission of ``The Optimizer Quotient and the Certification Trilemma'' for consideration for publication in the \textit{Journal of the ACM}.

\noindent\textbf{Summary}

This paper introduces the \emph{optimizer quotient}---the coarsest abstraction of a system's state that preserves optimal actions. We prove that certifying this object's coordinate structure is subject to an impossibility trilemma: under P$\neq$coNP, no certifier can be simultaneously sound, complete on all in-scope instances, and polynomial-budgeted. The certification cost varies by computational regime: coNP-complete (static), PP-hard (stochastic decisiveness), and PSPACE-complete (sequential). We identify six structural restrictions that collapse certification to polynomial time.

\noindent\textbf{Contributions}

\begin{itemize}
  \item \textbf{New complexity measure}: The optimizer quotient formalizes which state coordinates are decision-relevant.
  \item \textbf{Impossibility trilemma}: Under P$\neq$coNP, exact certification cannot be simultaneously sound, complete, and efficient.
  \item \textbf{Regime-aware complexity}: Certification hardness varies from coNP to PSPACE across static, stochastic, and sequential regimes.
  \item \textbf{Mechanized core}: The finite reduction and verification core is formalized in Lean 4.
\end{itemize}

\noindent\textbf{Significance}

This work contributes to the theory of computation by:
\begin{itemize}
  \item Establishing fundamental limits on automated decision certification
  \item Connecting decision-theoretic sufficiency to computational complexity
  \item Providing a framework for understanding when decision-relevant information can be efficiently certified
\end{itemize}

The results are of interest to the JACM audience in theoretical computer science, particularly in complexity theory, formal verification, and decision theory.

\noindent\textbf{Originality}

This work has not been published previously and is not under consideration for publication elsewhere. All authors have approved the submission.

\noindent\textbf{Suggested Reviewers}

We suggest the following reviewers who have expertise in relevant areas:

\begin{itemize}
  \item \textbf{[Reviewer 1]}: Expertise in computational complexity and proof complexity
  \item \textbf{[Reviewer 2]}: Expertise in formal verification and interactive proof systems
  \item \textbf{[Reviewer 3]}: Expertise in decision theory and information-theoretic complexity
\end{itemize}

\noindent\textbf{Conflict of Interest}

The authors declare no conflicts of interest.

Thank you for your time and consideration. We look forward to your response.

Sincerely,

\\
\textbf{Tristan Simas} \\
McGill University \\
Montreal, Quebec, Canada \\
\texttt{tristan.simas@mail.mcgill.ca}

\end{document}
